% !PRINTED 140
\weberpage{140}
%<<CONT
dass eine einfache Permutationsgruppe 60\textsuperscript{sten} Grades bei fünf Ziffern wirklich existirt, nämlich die alternirende Gruppe.

Man kann aber auch die Thatsache dadurch verificiren, dass man einen Theiler 12\textsuperscript{ten} Grades von $P$ nachweist (\S.\,28, 3.). Man erhält diesen Theiler 12\textsuperscript{ten} Grades z. B. daraus, dass die Gruppe 4\textsuperscript{ten} Grades:
\[
Q = 1, (1, 2) (3, 4), (1, 2) (5, 6), (3, 4) (5, 6)
\]
mit den Permutationen 3\textsuperscript{ten} Grades:
\[
c = (1, 3, 5) (2, 4, 6)
\]
vertauschbar ist, d. h. der Bedingung $c^{-1} Q c = Q$ genügt, und erhält also dann eine Gruppe 12\textsuperscript{ten} Grades:
\[
Q, \; Q c, \; Q c^2.
\]
Die einfache Gruppe 60\textsuperscript{sten} Grades wird auch die \gesperrt{Ikosaëdergruppe} genannt aus einem Grunde, den wir später kennen lernen werden.

Dieselbe Betrachtung zeigt auch, dass eine einfache Gruppe 90\textsuperscript{sten} Grades nicht existirt. Denn diese Gruppe müsste ein Element 5\textsuperscript{ten} und ein Element 3\textsuperscript{ten} Grades enthalten. Da $90 = 5 . 18$ ist und 18 keinen anderen Theiler als 6 hat, der nach dem Modul 5 mit 1 congruent ist, so muss, wenn wir für $Q$ eine Gruppe 5\textsuperscript{ten} Grades wählen, $R$ vom Index 6 sein und die Gruppe $P$ wäre also als transitive Permutationsgruppe von sechs Ziffern darstellbar. Ganz wie oben schliesst man, dass sie die beiden Elemente $a, b$ enthalten müsste, was, wie wir gesehen haben, zur Ikosaëdergruppe führt, die nicht Theiler einer Gruppe 90\textsuperscript{sten} Grades sein kann.

\begin{center}
\S.\,35.

\gesperrt{Gruppen vom Grade $pq$.}
\end{center}

Es ist nun noch von Interesse, zu untersuchen, wie bei einer gegebenen Gradzahl die Gruppen constituirt sind. Dafür sind die Sätze, die in den vorangegangenen Paragraphen abgeleitet sind, die Grundlagen. Freilich sind wir bis jetzt nur bei den einfacheren Gradzahlen im Stande, die vorhandenen Gruppen vollständig zu übersehen. Am weitesten geht hierin eine Arbeit von \gesperrt{Hölder}\footnote{\gesperrt{Hölder}, Die Gruppen der Ordnungen $p^3, pq^2, pqr, p^4$. Mathematische Annalen, Bd. 43.}. Wir betrachten hier nur den einfachen Fall,
%CONT>>
